% Generated draft: requires mathematical and source review.
\documentclass[11pt]{article}
\usepackage[margin=25mm]{geometry}
\usepackage{fontspec}
\setmainfont{Latin Modern Roman}
\usepackage{amsmath,amssymb,mathtools,unicode-math}
\setmathfont{Latin Modern Math}
\usepackage{enumitem,needspace,xurl,hyperref}
\setlength{\parindent}{0pt}
\setlength{\parskip}{5pt}
\setlength{\emergencystretch}{3em}
\newcommand{\N}{\mathbb{N}}
\newcommand{\Z}{\mathbb{Z}}
\newcommand{\Q}{\mathbb{Q}}
\newcommand{\R}{\mathbb{R}}
\newcommand{\C}{\mathbb{C}}
\newcommand{\F}{\mathbb{F}}
\newcommand{\A}{\mathbb{A}}
\newcommand{\PP}{\mathbb{P}}
\newcommand{\E}{\mathbb{E}}
\newcommand{\eps}{\varepsilon}
\newcommand{\dd}{\,\mathrm{d}}
\begin{document}
\section*{UnsolvedMath ID: 30002340}
\textbf{Source problem code:} catalogue entry
\subsection{Definitions and mathematical statement}
The retained source inquiry is
\begin{quote}
What is known about the Hecke algebra associated with the arithmetic group \(G(\Z)\), including its action on \(L^2(G(\Z)\backslash G(\R))\)?
\end{quote}
Here \(G\) is unspecified in the supplied record. If suitable integral and real points have been defined, write \(\Gamma=G(\Z)\); then \(\Gamma\backslash G(\R)\) denotes the left-coset space. With a specified quotient measure and scalar field, \(L^2(\Gamma\backslash G(\R))\) denotes the space of square-integrable functions, identified up to equality almost everywhere. The record does not specify these data or hypotheses ensuring the intended arithmetic and measure-theoretic setting.

The phrase ``the Hecke algebra associated with'' does not determine an algebra without further choices, such as an ambient group defining the relevant double cosets, or an adelic group and level subgroup. No such choices or precise action are supplied. The background discusses the discrete automorphic spectrum, whereas the retained inquiry displays the full \(L^2\)-space; the intended restriction is unresolved.

\textbf{Unresolved formulation qualification.} The supplied statement-verification field classifies this record as \texttt{unrecoverable}: the source's ``What does one know?'' is described as a rhetorical transition in an introductory account, answered by the following passage through the adelic framework. According to that supplied verification, no definite unresolved mathematical claim is stated, and no separate clean statement is promoted. Accordingly, the inquiry above is retained for provenance; it does not specify a conjecture, theorem to prove, or determinate classification problem. A precise mathematical target cannot be recovered from the supplied record without additional evidence.
\subsection{Short English statement}
The record asks broadly about Hecke algebras for \(G(\Z)\) and their action on \(L^2(G(\Z)\backslash G(\R))\). Its supplied verification identifies this as an expository transition rather than a definite open problem. The group, Hecke algebra, and intended spectral scope remain unspecified.
\subsection{Sources}
\begin{itemize}
\item Local UnsolvedMath registry: \texttt{data/unsolvedmath/problems.json}; permanent ID \(30002340\), exact source code \texttt{catalogue entry}, title ``Expository Transition to Adelic Hecke Operators.'' The supplied record belongs to the Oberwolfach Reports Open Problems collection. No individual proposer is supplied; the recorded proposed year is \(2013\).
\item Supplied report attribution: \emph{Mini-Workshop: Spherical Varieties and Automorphic Representations (2014)}, Oberwolfach Reports, DOI \texttt{10.4171/owr/2013/24}, pp.~1445--1494. \url{https://doi.org/10.4171/owr/2013/24}. The literature-source label is ``Oberwolfach Report 2013/24.'' The report is supplied as documentation of workshop discussions; the precise supporting passage is not included.
\item Supplied literature reference: Sakellaridis--Venkatesh, \emph{Periods and harmonic analysis on spherical varieties}. \url{https://arxiv.org/abs/1203.0039}.
\item Supplied literature assessment, checked on \texttt{2026-08-21}: the record labels the topic ``open'' and states that general Satake and spherical-variety theories provide important cases but no general description settles the extracted question. This is supplied dated curation, not an independently verified literature assessment. The background explicitly requests independent verification of time-sensitive status and cited claims.
\item Supplied statement verification: the inquiry is classified as \texttt{unrecoverable}, because the source question is described as an introductory rhetorical transition followed by an adelic explanation, with no definite unresolved claim. This qualification must be retained alongside the supplied ``open'' label.
\end{itemize}
\end{document}
